% This file is part of the Open Source project 'proTirone'
% (c) 2025 Karsten Reincke (https://github.com/protirone/protico.ltx)
% It is distributed under the terms of the creative commons license
% CC-BY-4.0 (= https://creativecommons.org/licenses/by/4.0/)

\selectlanguage{ngerman}


\begin{frame}  
  \frametitle{LFCX:BGE: \textcolor{red}{Müssen \& Sollen} $\rightarrow$ \textcolor{teal}{Dürfen \& Wollen}}


  \begin{columns}
    \column{0.4\textwidth}
      \begin{center}
      
\includegraphics[width=4cm]{\imgLf/control-885742-pxh.png}
      \end{center}
    \column{0.1\textwidth}
      \begin{center}
      $\rightarrow$
      \end{center}
    \column{0.4\textwidth}
      \begin{center}
      
\includegraphics[width=4cm]{\imgLf/trust-cat-dog-656221-pxh.png}
      \end{center}
    \end{columns}

\end{frame}


\begin{frame}  
  \frametitle{LFCX:BGE: Selbstverantwortung}

      \begin{center}
      
\includegraphics[height=7cm]{\imgLf/responsibility-1334149-pxh.png}
      \end{center}

\end{frame}

\begin{frame}  
  \frametitle{LFCX:BGE: Zuarbeit}

      \begin{center}
      \includegraphics[height=7cm]{\imgLf/support-1349598-pxh.png}
      \end{center}

\end{frame}

\begin{frame}  
  \frametitle{LFCX:BGE:Manifest für die GS-LDK}

  \begin{itemize}
    \item Wir halten Gesetze und Verordnungen ein:
      \begin{itemize}
        \footnotesize
        \item Schulpflicht, keine Gewalt, kein Mobbing, keine Störungen, ...
        \item Klausurteilnahme
        \item Entschuldigungen für Fehltage
        \item Beachtung der Höchstgrenze für Fehltage
        \item ...
      \end{itemize}
    \item Jede Schülerin erhält im Zeugnis bedingungslos eine 3.
    \item Klausurenergebnisse werden normiert auf die Noten 3 - 1.
    \item Für die mündliche Note sammeln die Schülerinnen Punkte:
      \begin{itemize}
        \footnotesize
      \item eine 3 : keine Punkte
      \item eine 2 : so viele Punkte, wie Unterrichtstage 
      \item eine 1 : 2,5 mal so viele Punkte, wie Unterrichtstage 
    \end{itemize}
    \end{itemize}

\end{frame}



\begin{frame}  
  \frametitle{LFCX:BGE:Möglichkeiten zum Punktesammln}

  \begin{itemize}
    \item Sehr aktive mündliche Mitarbeit: 1P/pro Schultag.
    \item Große Präsentation: 
      \begin{itemize}
      \footnotesize
      \item Unterlagen 1 Tag vorab: 1P
      \item Präsentieren: 1P
      \item inhaltliche Qualität, Nützlichkeit für Mitschüler, Präsentationsreife: je 1P
      \end{itemize}
    \item Große Hausarbeit: 1 - 3 Punkte nach Qualität
    \item Übungen (Diagramme etc.), Zusammenfassungen: 1P
    \item Supportive Behavior, besondere Unterstützung von Klassenkameradinnen: 2P
    \item Was mir ein unmittelbares *Wow* abnötigt: 1P
  \end{itemize}

\end{frame}

\begin{frame}  
  \frametitle{LFCX:BGE: Spirit}

      \begin{center}
      
\includegraphics[height=4cm]{\imgLf/im-grunde-gut.png}
      \end{center}

\end{frame}